% answers_short.tex —— 2016 年·春季上海卷 · 简短答案速查
% 由 tools/make_short_answers.py 自 answers.tex 的 \gksolhead/\gkanswer 提取，
% 并与 2016/2016_数学真题_春季上海_答案要点.md 逐题比对后生成，勿手改（改源 MD 后重跑）。

\begin{gknotice}
\textbf{使用说明}：本节只给最终答案便于快速核对，完整推导见解析版。答案为编写者独立推导，\textbf{非官方标准答案}。
\end{gknotice}

\gkpart{填空题}
\begin{itemize}[leftmargin=2.6em]
  \item[\textbf{1.}] \(3\)
  \item[\textbf{2.}] \(7\)
  \item[\textbf{3.}] \(\dfrac{\pi}{4}\)
  \item[\textbf{4.}] \([2,+\infty)\)
  \item[\textbf{5.}] \(8\)
  \item[\textbf{6.}] \(1\)
  \item[\textbf{7.}] \(2\sqrt{3}\)
  \item[\textbf{8.}] \(24\)
  \item[\textbf{9.}] \(3\)
  \item[\textbf{10.}] \(-4\)
  \item[\textbf{11.}] \([1,2]\)
  \item[\textbf{12.}] \(4\)
\end{itemize}

\gkpart{选择题}
\begin{center}\begin{tabular}{ccccccccccccc}
\toprule
题号 & 13 & 14 & 15 & 16 & 17 & 18 & 19 & 20 & 21 & 22 & 23 & 24 \\
\midrule
答案 & B & D & C & C & A & C & D & B & B & D & A & B \\
\bottomrule
\end{tabular}\end{center}

\gkpart{解答题}
\begin{itemize}[leftmargin=2.6em]
  \item[\textbf{25.}] \(\arccos\dfrac{3}{10}\)。
  \item[\textbf{26.}] 最小正周期 \(T=2\pi\)，最大值 \(2\)；取得最大值时 \(x=\dfrac{\pi}{6}+2k\pi\)（\(k\in\mathbf{Z}\)）。
  \item[\textbf{27.}] \(|OF|=3.6\ \mathrm{cm}\)。
  \item[\textbf{28.}] (1) \(a_{1}=-8\)；(2) \(c_{n}=n^{2}-20n+2^{n}-1\)，最小项 \(c_{n_{0}}=c_{4}=-49\)
  \item[\textbf{29.}] (1) \(D_{f>g}=(-\infty,-1)\cup(3,+\infty)\)；(2) \(a>-\dfrac{4}{9}\)
\end{itemize}

\gkpart{附加题}
\begin{itemize}[leftmargin=2.6em]
  \item[\textbf{30.}] B
  \item[\textbf{31.}] D
  \item[\textbf{32.}] B
  \item[\textbf{33.}] \(5\)
  \item[\textbf{34.}] \(50\pi\)
  \item[\textbf{35.}] \(28\)
  \item[\textbf{36.}] (1) \((1,3,5)\)，\((1,3,3)\)，\((1,2,5)\)，\((1,2,3)\)，共 4 个；(2) \(\{t\mid t\neq 3,\ t\neq 6,\ t\in\mathbf{N}^{\ast}\}\)
\end{itemize}
